% yoshi_fixed_v4.tex --- MIプログラム 中間発表予稿（A4・2段組・2ページ）
% コンパイル: platex yoshi_fixed_v4.tex && dvipdfmx yoshi_fixed_v4.dvi
\documentclass[a4paper,twocolumn,10pt]{jsarticle}
\usepackage{amsmath,amssymb}
\usepackage[dvipdfmx]{graphicx}
\usepackage{url}
\setlength{\columnsep}{7mm}
\setlength{\textheight}{264mm}
\setlength{\topmargin}{-21mm}
\setlength{\textwidth}{175mm}
\setlength{\oddsidemargin}{-7mm}
\setlength{\evensidemargin}{-7mm}
\setlength{\parindent}{1zw}
\pagestyle{empty}

\begin{document}

\twocolumn[{%
\begin{center}
{\normalsize 2026年10月2日(金)\quad MIプログラム\quad 中間発表}\par\vspace{2mm}
{\LARGE\bfseries 歪み関数重み付き1-Wasserstein距離に基づく}\par\vspace{1mm}
{\LARGE\bfseries quantile fraction学習とCVaR推定誤差の低減}\par\vspace{3mm}
{\normalsize 発表者: 【氏名】\qquad 所属: 【研究室名】\qquad 指導教員: 【指導教員名】}
\end{center}
\vspace{3mm}
}]

\section{はじめに}
金融派生商品のヘッジでは，損益分布の下側裾に生じる大きな損失をいかに抑えるかが本質的な問題となる．条件付きバリューアットリスク（CVaR）に基づくヘッジ方策の学習には，下側裾の正確な評価が不可欠である．分布型強化学習（分布型RL）はリターンの期待値ではなく分布そのものを学習する枠組みであり，リスク感応的な意思決定への応用が期待されている．しかし，QR-DQN，IQN\cite{iqn}，FQF\cite{fqf}といった分位点回帰型の既存手法は，真の分位点関数$F^{-1}$と近似分位点関数$\hat F^{-1}$の間の1-Wasserstein距離
\begin{equation}
W_1=\int_0^1 \bigl|F^{-1}-\hat F^{-1}\bigr|\,d\tau
\end{equation}
を実質的に最小化する．この距離は区間$[0,1]$のすべての$\tau$に一様な重みを与えるため，原子（離散分布に確率質量を置く点）の配置がリスク尺度と無関係な一様測度$d\tau$に対する近似誤差の評価に基づいて決まり，CVaRの評価に必要な下側裾の解像度が不十分となりやすい．FQFは配置を学習するが，目的関数が一様重みであるため，CVaR裾への集中を直接には促さない．ヘビーテイル分布（稀だが巨大な損失が無視できない型の分布）では，この解像度不足がCVaR推定誤差として顕在化する．本研究では，リスク尺度を規定する歪み関数を損失に組み込み，quantile fraction（累積確率軸 $[0, 1]$ の区切り位置）の配置自体をCVaR推定に対して最適化する手法を提案する．

\section{関連技術・先行研究}
分位点回帰型分布型RLは，リターン分布を有限個の分位点で表現し，分位点回帰損失によって学習する．FQF\cite{fqf}は図\ref{fig:fqf}に示すように，状態表現からquantile fraction（累積確率軸 $[0, 1]$ の区切り位置）$\tau_i$を出力するFraction Proposal Network（FPN）と，$\tau_i$をコサイン基底で埋め込んで分位点値を出力する分位点値ネットワークからなる．FPNは$W_1$を近似的に最小化するよう学習され，$\tau_i$は一様重みの下での最適量子化点を目指して学習される（収束の理論保証はFQFにも存在しない）．

\begin{figure}[t]
\centering
\includegraphics[width=0.75\columnwidth]{fig1_fqf_overview.png}
\caption{FQFの概要．状態表現からFPNがquantile fraction $\tau_i$を提案し，分位点値ネットワークが対応する分位点値を出力する．}
\label{fig:fqf}
\end{figure}

リスク尺度については，歪み関数$g$を用いた歪みリスク尺度が分位点関数の積分
\begin{equation}
\rho_g(Z)=\int_0^1 F^{-1}\,dg(\tau)
\end{equation}
で表され，$\mathrm{CVaR}_\alpha$は$g(\tau)=\min(\tau/\alpha,1)$に対応する\cite{rockafellar}．一方，1次元分布を有限個の点で近似する最適量子化の理論\cite{graf}は，重み付き$W_1$を最小化する原子配置の存在と収束率を与える．

先行研究として，IQN\cite{iqn}は歪みリスク尺度を$\tau$の再エンコードとして扱うのみで，分位点の配置は学習しない．Beyond CVaR\cite{beyondcvar}はスペクトルリスク尺度の重みを分位点の「値」に作用させるが，quantile fractionは一様に固定されている．Tail-Safe\cite{tailsafe}は$\tau$のサンプリングを下側裾へヒューリスティックに傾斜させる．いずれも，損失側でquantile fractionの配置をリスク尺度に対して最適化する定式化には至っていない．

\section{研究目的}
本研究の目的は，FPNの目的関数を一様重みの$W_1$から歪み関数重み付き$W_1$へ置換し，quantile fractionを歪み分布$\mathcal{D}_g Z$の1次元最適量子化点として学習する手法を確立することである．歪みリスク尺度の誤差は，既知のLipschitz評価の重み付き形\cite{pichler2013}として
\begin{equation}\small\begin{gathered}
|\rho_g(Z)-\rho_g(\hat Z)|\le\!\!\int_0^1\!|F^{-1}-\hat F^{-1}|g'(\tau)d\tau=W_1(\mathcal{D}_g Z,\mathcal{D}_g\hat Z)
\end{gathered}\normalsize\label{eq:bound}\end{equation}
と上から抑えられるため，この量を直接最小化することで，歪み重みに整合したCVaR推定誤差の低減を目指す．

\section{研究方針}
理論面では，歪み重み付き損失$L_g$のquantile fractionに関する勾配の閉形式を導出する．原子を$g$-中央値 $\hat\tau_i$ とすることで，
\begin{equation}
\frac{\partial L_g}{\partial\tau_i}=g'(\tau_i)\bigl[2F^{-1}(\tau_i)-F^{-1}(\hat\tau_i)-F^{-1}(\hat\tau_{i-1})\bigr]
\end{equation}
を得た．これを分位点値ネットワークにplug-inした擬似勾配更新（stop-gradient型）として実装する．数値面では，既知の分位点関数をもつ1次元分布での量子化デモ，自作toy環境での分布型RL比較，デルタヘッジ環境での実験の3段階で検証を進める．

\section{現在までの進捗}
まず新規性調査を完了し，前節の先行研究のいずれにも損失側でのquantile fraction最適化が存在しないことを確認した．

実験A（1次元量子化デモ）では，ガウス混合$0.9N(0,1)+0.1N(-4,1)$，負の対数正規，Student-t（$\nu=3$）の3分布に対し，$N\in\{8,\dots,256\}$，$\alpha\in\{0.01,0.05,0.1\}$の条件で，一様配置（M0），$g$-一様配置（M1），歪み最適配置（M2）のCVaR推定誤差を比較した．$\alpha N\ge 2$ 領域における最小改善比 $R=42.3$ を達成し，代表例 $\alpha=0.05, N=16$（一様配置のCVaR推定が単一分位点$\hat F^{-1}(1/32)$のみで行われる特殊設定）ではガウス混合で$R=97.9$，負の対数正規で$R=52.4$，Student-tで$R=39.6$となり，判定基準$R\ge 2$を満たした．歪み最適配置は15個の内部quantile fractionのうち13個を下側裾$[0,\alpha]$に配置し（一様は0個），点数換算値（N-equivalence）では歪み最適16点が一様695点に相当する（D1・$\alpha=0.05$での補間値；D2・D3は探索上限の1024点でも未達）．log-log傾きは歪み最適配置（M2）で$-1.79$〜$-1.84$，一様配置（M0）で$-0.72$〜$-1.27$であり，M2は二次減衰に，M0は一次減衰におおむね整合する（機序は下側裾の特異性に対する適応配置の回避．漸近次数の一般的証明は今後の課題）．改善比$R$は$N$とともに増大する（$N=1024$で$10^3$を超える）．図\ref{fig:demo}に$N$に対する誤差収束を示す．

\begin{figure}[t]
\centering
\includegraphics[width=0.72\columnwidth]{fig2_demo.png}
\caption{実験A: 原子数$N$に対するCVaR推定誤差のlog-log収束（行は$\alpha\in\{0.01,0.05,0.1\}$，列は3分布，本文中の数値は$\alpha=0.05$）．歪み最適配置（M2）は一様配置（M0）より一貫して小さい誤差を示す．}
\label{fig:demo}
\end{figure}

実験B（toy環境）では，ホライズン10，55状態$\times$2行動，$\gamma=1$でクラッシュによりヘビーテイル損失が生じる自作環境RiskyChain-v0を構築し，FQF（B0），歪みFQF（一様初期化B1および$g$-一様初期化B1g），リスク感応IQN（B2），学習時歪みサンプリングIQN（B2p）を100kステップ・3シードで学習した（CVaR-greedy，$\alpha=0.05$）．全110状態・行動対のCVaR推定MAE（mean$\pm$標本SD）は，B0が$7.02\pm0.14$，B1が$1.41\pm0.92$，主結果のB1gが$0.96\pm0.51$，B2が$0.36\pm0.05$，B2pが$3.42\pm1.38$であった．主結果B1gはFQF比で約1/7に誤差を低減し，3シード全てで理論配置（下側裾平均14.06個）を達成してB1の局所解を解消した．さらに，FPNを凍結して理論配置に固定したB3gは$0.74\pm0.33$であり，B1gとの差は3シードの標本SD内で優劣は確定できないが，toy環境では改善の主因がFPNの状態依存学習よりも理論配置の採用にあることが示唆された．なお，M-estの参照値は各手法のgreedy方策下の収益分布に対するものであり，手法間差には推定対象の難易度差が含まれる（共通方策下の交差評価でもall-RISKY参照の誤差は全手法で大きく，本評価は継続方策の不一致を含むため厳密な表現比較ではない）．方策評価（polCVaR）ではB1g=$0.70\pm0.43$，B2=$1.01\pm0.01$であり，閾値方策クラス内最良（$k^*=0$）の値$1.0000$（収益確定）に近い水準に達した（B0は$-10.79\pm0.00$，評価用乱数シード固定のためSD$=0.00$）．一方，B2pは学習時サンプリングを歪めたことで約10倍悪化（$0.36\to3.42$）した（ターゲット側歪みサンプリングによる系統的バイアスの解釈があり，sampling系一般の劣位を示すものではない）．ただしIQN（B2）が全体として最小誤差を示しており，一般的優位性は未確認である．計算はRyzen 7 5825U・32GB RAMのCPU演算のみで行った．

\section{今後の予定}
10月に金融デルタヘッジ環境の実装，11月に本実験とベースライン比較，12月に理論証明の整備と論文執筆，1月に最終発表準備を行う．

\begin{thebibliography}{9}
\scriptsize
\setlength{\itemsep}{-2pt}
\setlength{\parskip}{0pt}
\bibitem{fqf} D. Yang ほか. Fully Parameterized Quantile Function for Distributional Reinforcement Learning. NeurIPS 2019.
\bibitem{iqn} W. Dabney ほか. Implicit Quantile Networks for Distributional Reinforcement Learning. ICML 2018.
\bibitem{beyondcvar} M. Moghimi, H. Ku. Beyond CVaR: Leveraging Static Spectral Risk Measures for Enhanced Decision-Making in Distributional Reinforcement Learning. ICML 2025.
\bibitem{tailsafe} J. Zhang. Tail-Safe Hedging. arXiv:2510.04555, 2025.
\bibitem{graf} S. Graf, H. Luschgy. Foundations of Quantization for Probability Distributions. LNM 1730, Springer, 2000.
\bibitem{rockafellar} R. T. Rockafellar, S. Uryasev. Optimization of Conditional Value-at-Risk. Journal of Risk 2(3), 2000.
\bibitem{pichler2013} A. Pichler. Evaluations of Risk Measures for Different Probability Measures. SIAM J. Optim. 23(1), 2013.
\end{thebibliography}

\end{document}
